% Study architecture (hand-authored TikZ): an illustrated left->right pipeline
% (data heatmap -> two parallel FM/baseline lanes with a transformer stack and UMAP
% scatters -> a fair-evaluation module -> three finding mini-plots), not labelled boxes.
{\hyphenpenalty=10000\exhyphenpenalty=10000\relax%
\definecolor{cUMa}{HTML}{1b7837}\definecolor{cUMb}{HTML}{762a83}\definecolor{cUMc}{HTML}{2c7fb8}%
% --- reusable pictograms -------------------------------------------------------
% gene x cell expression heatmap (rows=genes, cols=cells)
% NOTE: every macro-definition line below ends in "%" -- without it, the trailing
% newline becomes a real (invisible, inkless) interword space appended to the
% horizontal list \resizebox measures, inflating this file's natural width beyond
% what TikZ's own "current bounding box" tracks and de-centring the panel when
% resized (9 unprotected lines here were worth ~30pt of hidden slack).
\newcommand\heatmap[1]{\begin{tikzpicture}[scale=#1,baseline=(current bounding box.center)]
  \foreach \c/\r/\v in {0/0/25,0/1/60,0/2/10,1/0/70,1/1/15,1/2/45,2/0/20,2/1/80,2/2/35,
                        3/0/55,3/1/30,3/2/65,4/0/40,4/1/50,4/2/20}{
    \fill[cData!\v] (\c*2.6pt,\r*2.6pt) rectangle ++(2.4pt,2.4pt);}
  \draw[cData!60,line width=0.3pt] (-0.6pt,-0.6pt) rectangle (13.4pt,8.2pt);
\end{tikzpicture}}%
% transformer block stack (foundation model)
\newcommand\transformer{\begin{tikzpicture}[baseline=(current bounding box.center)]
  \foreach \y in {0,1,2,3}{\draw[cFM!70,line width=0.5pt,fill=cFM!14,rounded corners=1pt]
    (0,\y*3.3pt) rectangle ++(17pt,2.5pt);}
\end{tikzpicture}}%
% small UMAP scatter with three colour clusters
\newcommand\umap[1]{\begin{tikzpicture}[scale=#1,baseline=(current bounding box.center)]
  \foreach \x/\y in {0/0,1.4/0.6,0.6/1.5,-0.8/0.9,0.3/-0.9}{\fill[cUMa] (\x pt,\y pt) circle(0.9pt);}
  \foreach \x/\y in {7/5,8.2/5.8,6.4/6.2,7.6/4.1,8.8/5.1}{\fill[cUMb] (\x pt,\y pt) circle(0.9pt);}
  \foreach \x/\y in {4/-3,5.2/-2.2,3.4/-2,4.6/-3.9,5.6/-3.1}{\fill[cUMc] (\x pt,\y pt) circle(0.9pt);}
\end{tikzpicture}}%
% balance-scale glyph (matched / fair comparison)
\newcommand\scaleicon{\begin{tikzpicture}[baseline=(current bounding box.center)]
  \draw[cEval,line width=0.8pt] (0,6pt)--(0,-1pt); \draw[cEval,line width=0.8pt] (-5pt,6pt)--(5pt,6pt);
  \draw[cEval,line width=0.6pt] (-5pt,6pt)--(-6.5pt,1.5pt)--(-3.5pt,1.5pt)--cycle;
  \draw[cEval,line width=0.6pt] (5pt,6pt)--(3.5pt,1.5pt)--(6.5pt,1.5pt)--cycle;
  \draw[cEval,line width=0.8pt] (-2.5pt,-1pt)--(2.5pt,-1pt);
\end{tikzpicture}}%
% --- three finding mini-plots (stylised, matching the real result shapes) ------
% F1 parity: FM vs PCA on the diagonal
\newcommand\miniparity{\begin{tikzpicture}[baseline=(current bounding box.center)]
  \draw[black!35,line width=0.4pt] (0,0) rectangle (17pt,17pt);
  \draw[black!30,dashed,line width=0.4pt] (1pt,1pt)--(16pt,16pt);
  \foreach \x/\y in {5/6,8/7.5,10/11,12/12.5,7/8,13/12,9/9.5,11/10.5}{\fill[cFind] (\x pt,\y pt) circle(0.7pt);}
\end{tikzpicture}}%
% F2 collapse: scRNA rises with shift, scATAC flat
\newcommand\minicollapse{\begin{tikzpicture}[baseline=(current bounding box.center)]
  \draw[black!35,line width=0.4pt] (0,0) rectangle (17pt,17pt);
  \draw[cUMc,line width=0.9pt] (1pt,2pt) .. controls (9pt,3pt) and (12pt,10pt) .. (16pt,15pt);
  \draw[cData,line width=0.9pt] (1pt,3pt)--(16pt,3.6pt);
\end{tikzpicture}}%
% F3 not-scale: flat across model size
\newcommand\miniscale{\begin{tikzpicture}[baseline=(current bounding box.center)]
  \draw[black!35,line width=0.4pt] (0,0) rectangle (17pt,17pt);
  \foreach \x in {3,8,13}{\fill[cFM!70] (\x pt-1.6pt,1pt) rectangle ++(3.2pt,10.5pt);}
  \draw[black!45,dashed,line width=0.4pt] (1pt,11.5pt)--(16pt,11.5pt);
\end{tikzpicture}}%
% ------------------------------------------------------------------------------
\begin{tikzpicture}[font=\footnotesize,>=Stealth,
  stg/.style={font=\scriptsize\bfseries,text=#1},
  sub/.style={font=\tiny,text=black!60,align=center},
  lane/.style={rounded corners=3pt,draw=#1!55,fill=#1!6,inner sep=3pt},
  flow/.style={-{Stealth[length=4pt,width=4pt]},line width=1pt,draw=black!35}]
% --- Stage 1: DATA -----------------------------------------------------------
\node (hm) at (0,0.55) {\heatmap{1.2}};
\node[stg=cData,above=2pt of hm] {DATA};
\node[sub,below=2pt of hm,text width=2.0cm] {scRNA $\cdot$ spatial \\ scATAC $\cdot$ perturb\\[1pt]\emph{$+12$ external}};
% --- Stage 2: two embedding lanes (FM top, baseline bottom) ------------------
% Horizontal stage spacing compressed and the two lanes pushed further apart vertically
% so the panel reads as a denser block (less wide-thin dead space) closer to panel B's aspect.
\node[lane=cFM,minimum width=2.6cm,minimum height=0.7cm] (fm) at (3.05,1.4) {\transformer};
\node[stg=cFM,above=1pt of fm] {FOUNDATION MODEL};
\node[lane=cBase,minimum width=2.6cm,minimum height=0.7cm] (bl) at (3.05,-0.3)
  {\footnotesize\color{cBase!85}HVG\,$\to$\,PCA\,/\,scVI};
\node[stg=cBase,below=1pt of bl] {SIMPLE BASELINE};
% --- Stage 3: one UMAP embedding per lane -----------------------------------
\node[draw=black!30,rounded corners=2pt,minimum size=0.7cm] (fu) at (5.45,1.4) {\umap{0.9}};
\node[draw=black!30,rounded corners=2pt,minimum size=0.7cm] (bu) at (5.45,-0.3) {\umap{0.9}};
\node[sub] at (5.45,0.55) {cell\\embeddings};
% --- Stage 4: fair evaluation -----------------------------------------------
\node (sc) at (7.65,0.55) {\scaleicon};
\node[stg=cEval,above=6pt of sc] {FAIR EVAL};
\node[sub,below=1pt of sc,text width=2.0cm] {$k$NN $\cdot$ ref-free $R^2$\\ conformal $\cdot$ $\pm0.02$};
% --- Stage 5: three findings (stylised mini-plots) --------------------------
\node (g1) at (9.75,0.72) {\miniparity};
\node[right=4pt of g1] (g2) {\minicollapse};
\node[right=4pt of g2] (g3) {\miniscale};
\node[draw=black!22,rounded corners=2.5pt,fit=(g1)(g2)(g3),inner sep=2.5pt] (gbox) {};
\node[stg=cFind] at ($(g2)+(0,0.62)$) {RESULTS};
\node[sub,below=1pt of g1,text width=0.95cm] {parity};
\node[sub,below=1pt of g2,text width=0.95cm] {collapse\\is general};
\node[sub,below=1pt of g3,text width=0.95cm] {not scale};
% --- flow arrows -------------------------------------------------------------
\draw[flow] (hm.east) to[out=30,in=180] (fm.west);
\draw[flow] (hm.east) to[out=-30,in=180] (bl.west);
\draw[flow] (fm.east) -- (fu.west);
\draw[flow] (bl.east) -- (bu.west);
\draw[flow] (fu.east) to[out=-22,in=150] (sc.west);
\draw[flow] (bu.east) to[out=22,in=210] (sc.west);
\draw[flow] (sc.east) to[out=0,in=180] (gbox.west);
% panel label at the picture's true top-left. Both Fig-1 panels are now \resizebox'd to the
% SAME width (0.88\textwidth) and \centering'd, so their bounding boxes share a left edge; with
% xshift=0 here and in fig0b_roadmap.tex the A and B tags land on one vertical column.
\node[anchor=south west,font=\normalsize\bfseries] at ([xshift=0pt,yshift=2pt]current bounding box.north west) {A};
\end{tikzpicture}}
